\section{问题二的建立与求解}

\subsection{数据驱动的BMI分组}

\subsubsection{选择依据}

不同BMI水平的孕妇，胎儿游离DNA浓度的增长速率存在差异，NIPT检测达标时间随BMI变化。若对全部孕妇采用统一检测时点，高BMI群体将面临较高的误检风险，需要按BMI水平分组并分别确定最优检测时点。传统做法依赖临床经验设定固定分界值，缺乏统计支撑。本文采用数据驱动策略，以组内达标时间方差最小化为目标，自动搜索最优分组断点。

\subsubsection{总述}

267名孕妇中，260名在观测期内Y染色体浓度首次达到检测阈值$c=0.04$（即4\%），据此定义经验达标时间$T^*$。对$T^*$关于BMI建立\textbf{Huber稳健回归}模型，在候选断点空间中穷举搜索，以BIC准则同时确定分组数目与断点位置。

\subsubsection{求解步骤}

\textbf{Step 1\quad 经验达标时间提取。}对每位孕妇的纵向浓度序列，取Y染色体浓度首次满足$\hat{y}(t) \geq 0.04$的孕周作为经验达标时间$T^*$。267名孕妇中有260名在观测窗口内达标，后续分析基于这260条记录。

\textbf{Step 2\quad 稳健回归建模。}以BMI为自变量，$T^*$为因变量，拟合二次稳健回归模型
\begin{equation}\label{eq:robust_reg}
  T^* = \beta_0 + \beta_1 \cdot \mathrm{BMI} + \beta_2 \cdot \mathrm{BMI}^2,
\end{equation}
采用Huber M-估计量以降低离群点的影响。图~\ref{fig:threshold_bmi}展示了$T^*$随BMI的变化趋势与拟合曲线。

\textbf{Step 3\quad 穷举断点搜索。}在BMI$\in[24,44]$范围内以步长1.0枚举所有候选断点组合，约束每组BMI宽度不少于3.0个单位、每组样本量不少于5人。对每种分组方案，计算组内$T^*$的加权方差，并以BIC准则评价模型复杂度与拟合优度的平衡。

\textbf{Step 4\quad 最优分组确定。}BIC最小方案对应两个断点BMI$=30$与BMI$=36$，将孕妇分为3组。临床经验5组方案（断点20/28/32/36/40）的BIC更高，额外分组并未带来组内方差的实质性下降。

\begin{figure}[H]
  \centering
  \includegraphics[width=0.85\textwidth]{Q2/experiments/round1/figures/M1_threshold_time_vs_bmi.png}
  \caption{经验达标时间$T^*$与BMI的关系及稳健回归拟合曲线，竖线标注最优断点}\label{fig:threshold_bmi}
\end{figure}

\subsection{风险函数的建立}

\subsubsection{选择依据}

NIPT检测时点的选择涉及两类竞争风险。检测过早，胎儿游离DNA浓度尚未充分积累，准确性不足；检测过晚，则延误临床干预窗口。需要构造统一的风险函数，将两类风险量化并加权求和，为最优时点求解提供目标函数。

\subsubsection{总述}

定义总风险函数为延迟风险与不准确风险的加权和，其中不准确风险基于问题一的GAMM预测值与残差标准差，通过正态累积分布函数刻画浓度低于阈值的概率。

\subsubsection{模型构建}

总风险函数定义为
\begin{equation}\label{eq:total_risk}
  R(t,\mathrm{BMI}) = w_d \cdot R_{\text{delay}}(t) + w_a \cdot R_{\text{inacc}}(t,\mathrm{BMI}),
\end{equation}
其中$w_d=0.4$，$w_a=0.6$，反映临床实践中检测准确性的优先级高于时效性。

延迟风险采用Sigmoid函数建模
\begin{equation}\label{eq:delay_risk}
  R_{\text{delay}}(t) = \frac{1}{1+e^{-0.5(t-18)}},
\end{equation}
该函数在孕18周附近快速上升，刻画超过常规检测窗口后临床干预机会急剧减少的特征。

不准确风险基于GAMM预测浓度$\hat{y}(t,\mathrm{BMI})$与残差标准差$\sigma_\varepsilon=0.0101$，通过标准正态累积分布函数$\Phi$计算浓度低于检测阈值$c=0.04$的概率
\begin{equation}\label{eq:inacc_risk}
  R_{\text{inacc}}(t,\mathrm{BMI}) = \Phi\!\left(\frac{c - \hat{y}(t,\mathrm{BMI})}{\sigma_\varepsilon}\right).
\end{equation}

当预测浓度远高于阈值时，$R_{\text{inacc}}$趋近于0；当预测浓度接近或低于阈值时，$R_{\text{inacc}}$迅速增大。图~\ref{fig:risk_analysis}展示了两类风险分量及总风险随孕周的变化。

\begin{figure}[H]
  \centering
  \includegraphics[width=0.85\textwidth]{Q2/experiments/round1/figures/M1_risk_function_analysis.png}
  \caption{各BMI组的风险函数分析，延迟风险、不准确风险及总风险随孕周的变化}\label{fig:risk_analysis}
\end{figure}

\subsection{最优检测时点求解}

\subsubsection{选择依据}

风险函数$R(t,\mathrm{BMI})$在孕周$t$上呈先降后升的"U形"结构，早期以不准确风险为主导，晚期以延迟风险为主导。两类风险的交汇区域对应总风险最低的检测窗口，该最低点即为各组最优检测时点。

\subsubsection{求解方法}

对第$g$组，取该组BMI均值$\overline{\mathrm{BMI}}_g$代入风险函数，在孕周区间$[10,25]$上求解单变量有界优化问题
\begin{equation}\label{eq:opt_time}
  T_g^* = \arg\min_{t\in[10,25]} R(t,\,\overline{\mathrm{BMI}}_g),
\end{equation}
使用Brent方法（scipy.optimize.minimize\_scalar）进行精确求解。

\subsubsection{求解结果}

三组的最优检测时点及对应风险值列于表~\ref{tab:opt_results}。

\begin{table}[H]
  \centering
  \caption{各BMI组最优检测时点与风险值}\label{tab:opt_results}
  \begin{tabular}{cccccc}
    \toprule
    组别 & BMI范围 & 样本量 & BMI均值 & 最优孕周 & 最小风险 \\
    \midrule
    G1 & $[20.7,\,30.0)$ & 76  & 28.9 & 11.72 & 0.0208 \\
    G2 & $[30.0,\,36.0)$ & 171 & 32.3 & 11.80 & 0.0218 \\
    G3 & $[36.0,\,46.9)$ & 20  & 38.4 & 13.49 & 0.0675 \\
    \bottomrule
  \end{tabular}
\end{table}

G1与G2组的最优时点十分接近（约11.7$\sim$11.8周），最小风险均低于0.025，中低BMI孕妇在孕12周前即可获得高置信度的检测结果。G3组（BMI$\geq$36）的最优时点推迟至13.49周，最小风险升高至0.0675，约为前两组的3倍。高BMI孕妇的母体游离DNA对胎儿DNA的稀释效应更强，浓度达标时间整体后移。

经验达标时间的中位数（G1为12.79周、G2为12.71周、G3为16.00周）均晚于模型给出的最优时点，基于风险函数的最优解能够在保证准确性的前提下提前检测窗口。

\subsection{检测误差敏感性分析}

\subsubsection{分析目的}

实际检测中，仪器精度与样本处理过程均可能引入系统偏差或随机噪声。需要评估检测阈值偏移或残差方差增大时，各组最优时点与最小风险的稳定性。

\subsubsection{扰动设计}

设置两类扰动实验。

（1）阈值偏移。将检测阈值$c$分别上调0.5\%、1\%、2\%（即$c=0.045$、$0.05$、$0.06$），模拟系统性测量偏差。

（2）噪声膨胀。将残差标准差$\sigma_\varepsilon$分别放大至1.5倍、2倍、3倍（即$\sigma_\varepsilon=0.0152$、$0.0202$、$0.0303$），模拟随机误差增大的场景。

\subsubsection{敏感性结果}

图~\ref{fig:sensitivity}展示了各组在不同扰动水平下的风险响应。G1和G2组在所有扰动条件下，最优时点偏移不超过0.5周，最小风险增幅有限，鲁棒性较好。G3组对两类扰动均敏感，噪声放大至3倍时最小风险超过0.15，最优时点进一步推迟。这与高BMI组信噪比偏低的物理机制一致，浓度预测的不确定性在高BMI区间被放大。

\begin{figure}[H]
  \centering
  \includegraphics[width=0.85\textwidth]{Q2/experiments/round1/figures/M1_error_sensitivity.png}
  \caption{检测误差敏感性分析，阈值偏移与噪声膨胀对各组风险的影响}\label{fig:sensitivity}
\end{figure}

作为对照，临床经验5组方案（M2，断点为BMI$=$20/28/32/36/40）给出的最优检测时点分布在11.65$\sim$14.65周之间，趋势与数据驱动方案一致，但额外的分组未能带来风险的明显降低，验证了BIC准则选择3组方案的合理性。

\subsection{实际意义}

上述结果为临床NIPT检测方案的制定提供了量化参考。BMI低于30的孕妇（占样本的28.5\%）和BMI在30$\sim$36之间的孕妇（占64.0\%）可在孕11$\sim$12周安全接受检测，风险水平低于0.025。BMI超过36的孕妇（占7.5\%）应将检测时点推迟至约13.5周，临床解读时需考虑该群体面临的较高检测不确定性。对于高BMI孕妇，在条件允许时可考虑重复采样或联合其他产前筛查手段，以弥补单次NIPT信噪比不足的问题。
